\documentclass[11pt,a4paper]{article}
\usepackage[margin=25mm]{geometry}
\usepackage{amsmath,amssymb,mathtools,enumitem}
% xurl is absent from this machine's TinyTeX; url + UrlBreaks gives the same
% break-anywhere behaviour for the long file paths cited in \source{}.
\usepackage{url}
\makeatletter\g@addto@macro\UrlBreaks{\UrlOrds}\makeatother
\usepackage[hidelinks]{hyperref}
\renewcommand{\texttt}[1]{\nolinkurl{#1}}
\setlist[itemize]{nosep,leftmargin=*}
\newcommand{\trans}{\mathsf{T}}
\newcommand{\source}[1]{\par\smallskip\noindent\textit{Source: #1.}\par\smallskip}
\title{EE6221 Robotics and Intelligent Sensors\\\large Beginner Notes: Manipulator Kinematics}
\author{Living course note}
\date{Updated: 24 August 2026}
\begin{document}
\maketitle
\noindent\textbf{Scope and sources.} This self-contained note develops the supplied manipulator material from first vocabulary through direct and inverse kinematics, elementary trajectory planning, Jacobians, and singularities. It uses \texttt{resources/lecture-01-manipulator-kinematics.pdf}, pp.~1--156, \texttt{resources/course-outline-part-01.pdf}, and the Week~1/2 transcripts as supporting teaching context. Newly archived public slides, assignments, and examinations are recorded in \texttt{STATUS.md} but are not treated as covered here until independently reviewed.
\tableofcontents

\section{\enspace What is a robot manipulator?}
For this course, a robot is an industrial, reprogrammable, multifunctional manipulator. A manipulator is a chain of rigid \emph{links} joined by actuated \emph{joints}; the final link carries an \emph{end effector}, such as a gripper, hand, or tool. Moving the joints changes the end effector's \emph{pose}: its position and orientation.

The two fundamental joint types are:
\begin{itemize}
\item \textbf{Revolute (R):} the link rotates about an axis; its joint coordinate is an angle.
\item \textbf{Prismatic (P):} the link slides along an axis; its joint coordinate is a displacement.
\end{itemize}
The number of independently controlled joint coordinates is the robot's degrees of freedom (DoF). The first few joints often position the wrist, while wrist joints orient the tool. A robot can be useful with fewer than six DoF, but then it may be unable to realise arbitrary position-and-orientation requests.

\section{\enspace Why and how robots are used}
Robots are particularly useful for repetitive, hazardous, precise, or reconfigurable tasks. In contrast to a hard-automation machine designed for one fixed task, a reprogrammable robot can change task by changing its program and tooling. This flexibility has a cost: it may be less economical than dedicated equipment at very high production volumes, while manual work can be preferable at very low volumes.

Common manipulator structures are named from the major joint types: Cartesian PPP, cylindrical RPP, spherical RRP, SCARA RRP, and articulated RRR. Their structure affects the workspace: the set of tool poses that can be reached. The workspace is not simply a ball around the base. Joint limits, link dimensions, collisions, and orientation requirements restrict it.

Useful specifications include payload (allowable load), reach, stroke, speed, repeatability, accuracy, and resolution. \emph{Accuracy} is closeness to the commanded pose; \emph{repeatability} is how consistently the robot returns to a pose. A robot may be repeatable but inaccurate if a systematic calibration error shifts every placement.
\source{\texttt{resources/lecture-01-manipulator-kinematics.pdf}, pp.~1--22; \texttt{resources/week-01-transcript.txt}.}

\section{\enspace Why coordinate frames are essential}
Saying ``the tool is at $(1,2,3)$'' is incomplete: those numbers only have meaning relative to a coordinate frame. We write $\{A\}$ for a base or reference frame and $\{B\}$ for a frame attached to a robot link. A physical vector is the same geometric object, but its components change when expressed in another frame.

If ${}^{A}\!R_B$ describes frame $\{B\}$ relative to frame $\{A\}$, then
\[
 {}^{A}\!v={}^{A}\!R_B\,{}^{B}\!v.
\]
The columns of ${}^{A}\!R_B$ are the $x_B$, $y_B$, and $z_B$ axes written in frame $\{A\}$. This gives an immediate interpretation: a rotation matrix tells an observer in $\{A\}$ which way the axes attached to $\{B\}$ point.

\section{\enspace Rotations and homogeneous transformations}
A valid three-dimensional rotation matrix satisfies
\[
 R\trans R=I,\qquad R^{-1}=R\trans,\qquad\det R=+1.
\]
For example, a right-handed rotation about the $z$ axis is
\[
 R_z(\theta)=\begin{bmatrix}
 \cos\theta&-\sin\theta&0\\
 \sin\theta&\cos\theta&0\\
 0&0&1
 \end{bmatrix}.
\]
Matrix order matters. Rotations about fixed axes premultiply an existing orientation; rotations about moving (body) axes postmultiply. Always state the frame and physical order instead of relying on a memorised product.

Rotation alone cannot account for different origins. If the origin of $\{B\}$ is at ${}^{A}\!p_B$, a point $P$ satisfies
\[
 {}^{A}\!p_P={}^A\!R_B{}^B\!p_P+{}^A\!p_B.
\]
Homogeneous coordinates combine both operations:
\[
 {}^A\!T_B=\begin{bmatrix}{}^A\!R_B&{}^A\!p_B\\0&1\end{bmatrix},\qquad
 \begin{bmatrix}{}^A\!p_P\\1\end{bmatrix}={}^A\!T_B\begin{bmatrix}{}^B\!p_P\\1\end{bmatrix}.
\]
The inverse is
\[
 ({}^A\!T_B)^{-1}=\begin{bmatrix}({}^A\!R_B)\trans&-({}^A\!R_B)\trans{}^A\!p_B\\0&1\end{bmatrix}.
\]
The translated term must be rotated when inverting; simply negating it is a common mistake.
\source{\texttt{resources/lecture-01-manipulator-kinematics.pdf}, pp.~23--52; \texttt{resources/week-01-transcript.txt}.}

\section{\enspace Direct kinematics and D--H construction}
Direct kinematics asks: if the joint coordinates $q_1,\ldots,q_n$ are known, where is the tool? The answer is the product
\[
 {}^0T_n(q)={}^0T_1(q_1){}^1T_2(q_2)\cdots{}^{n-1}T_n(q_n).
\]
The superscripts and subscripts behave like frame labels: adjacent labels cancel in a valid chain. If a product does not follow a connected route from the source frame to the destination frame, its physical meaning is suspect.

The standard Denavit--Hartenberg construction represents each adjacent-link transform by four operations:
\[
 {}^{i-1}T_i=R_z(\theta_i)\,T_z(d_i)\,T_x(a_i)\,R_x(\alpha_i).
\]
Here $\theta_i$ rotates about $z_{i-1}$, $d_i$ translates along $z_{i-1}$, $a_i$ translates along $x_i$, and $\alpha_i$ rotates about $x_i$. For a revolute joint, $\theta_i$ is the variable; for a prismatic joint, $d_i$ is the variable. The other quantities describe fixed link geometry.

Use this construction order:
\begin{enumerate}
\item Draw every joint axis and assign it as a $z$ axis.
\item Choose each $x_i$ along the common normal from $z_{i-1}$ to $z_i$; intersecting axes give $a_i=0$.
\item Complete a right-handed frame with $y_i=z_i\times x_i$.
\item Fill the D--H table before multiplying matrices.
\item Check the final rotation block is orthonormal and the translation has length units.
\end{enumerate}

\paragraph{Where each frame's origin actually lands.}
The two translation parameters do not just appear in the matrix --- they \emph{move the origin}, and
knowing where each origin ends up is what makes the table readable. Reading
${}^{k-1}T_k=\mathrm{Rot}(z,\theta_k)\mathrm{Trans}(z,d_k)\mathrm{Trans}(x,a_k)\mathrm{Rot}(x,\alpha_k)$
left to right: $d_k$ slides the origin along $z_{k-1}$, then $a_k$ slides it along the new $x_k$. So
the translation column of ${}^{k-1}T_k$ is $(a_kc_k,\ a_ks_k,\ d_k)$, and frame $k$'s origin sits
$a_k$ away from frame $k-1$'s whenever $a_k\neq0$.

For the articulated (RRR) waist/shoulder/elbow arm, with rows
$(\theta_1,d_1,0,90^\circ)$, $(\theta_2,0,a_2,0)$, $(\theta_3,0,a_3,0)$, the origins therefore land
on the physical landmarks:
\begin{itemize}
\item frame 1 --- the \textbf{shoulder}, a height $d_1$ up the column ($a_1=0$, so no sideways step);
\item frame 2 --- the \textbf{elbow}, a distance $a_2$ out along link 2;
\item frame 3 --- the \textbf{wrist}, a distance $a_3$ out along link 3.
\end{itemize}
This is exactly why $a_2$ and $a_3$ \emph{are} the link lengths: $x_k$ runs along the link, so the
step along $x_k$ is the link.

\paragraph{Common pitfall.}
$a_k=0$ does not mean frames $k-1$ and $k$ coincide. It means frame $k$'s origin stays
\emph{on the $z_{k-1}$ axis} --- there is no common normal to step along, because the two joint axes
intersect. The frames can still be separated by $d_k$ along that axis, and they are still different
frames with different orientations. In the arm above, $a_1=0$ puts the shoulder origin on the column
axis; frames 1 and 2 remain a full link apart.

\paragraph{One row per joint --- and a rigid tool is not a joint.}
The table has exactly as many rows as the arm has \emph{joints}. It is tempting to give a row to
every offset you can see in the drawing, but a fixed tool of length $L_t$ bolted to the last link
contributes no degree of freedom, so it must never occupy a joint row. Handle it in one of two ways,
and say which you are doing:
\begin{itemize}
\item post-multiply the arm matrix by a constant $\mathrm{Trans}(z,L_t)$, giving
      ${}^0T_{\text{tool}}={}^0T_n\,\mathrm{Trans}(z,L_t)$; or
\item add an explicitly labelled \emph{fixed} frame after the last joint row, with no starred
      variable in it.
\end{itemize}
If a tool row displaces a real joint, the chain silently loses a degree of freedom: the forward
equations can no longer reach the poses the arm physically reaches, and the inverse solution will
ask for a joint variable that the table never defined.

\paragraph{The thirty-second pre-flight check.}
Before you multiply anything, count. Mark each joint variable in the table with a star
($\theta_k^\ast$ or $d_k^\ast$), then count the stars and count the joints in the figure. The two
numbers must agree. A cylindrical arm drawn with one revolute and two prismatic joints needs three
starred variables --- one $\theta$ and two $d$'s. If you have two stars and a constant $L_t$ sitting
in a row, a joint has gone missing. This check costs seconds and catches the single most expensive
error in the question, because every later part inherits it.

A second consistency check is just as cheap: the $\alpha_k$ in your table and the
$R_x(\alpha_k)$ you write in ${}^{k-1}T_k$ must carry \emph{the same sign}. Writing
$\alpha_2=+90^\circ$ in the table and $R_x(-90^\circ)$ in the transform is worse than either
error alone, because a marker cannot tell which one you meant and cannot award the benefit of the
doubt.

For a planar two-revolute-link arm,
\[
 x=a_1\cos q_1+a_2\cos(q_1+q_2),\qquad
 y=a_1\sin q_1+a_2\sin(q_1+q_2).
\]
The second link uses the accumulated angle $q_1+q_2$ because its orientation is measured after both joint rotations.
\source{\texttt{resources/lecture-01-manipulator-kinematics.pdf}, pp.~53--86.}

\section{\enspace Inverse kinematics: pose to joint coordinates}
Inverse kinematics reverses the question: given a desired tool transform $^0T_n^*$, find $q$ such that $^0T_n(q)= {}^0T_n^*$. This is not ordinary matrix inversion because $q$ appears inside trigonometric functions and the mapping can be many-to-one.

Before solving equations, check three forms of feasibility:
\begin{itemize}
\item the desired position lies in the work envelope;
\item the robot has enough orientation freedom for the desired tool orientation;
\item each candidate respects joint limits and practical collision constraints.
\end{itemize}
An IK problem can have no solution, one solution, or multiple branches such as elbow-up and elbow-down. Use $\operatorname{atan2}(y,x)$ when reconstructing an angle: unlike $\arctan(y/x)$, it preserves quadrant information and remains meaningful when $x=0$.

\paragraph{Finish the answer: feasibility and how many solutions.}
Solving for the joint variables is only two-thirds of an inverse-kinematics answer. A complete
solution also states the \emph{range} over which it is valid and \emph{how many} solutions there
are. For a cylindrical arm with tool-tip target $(x,y,z)$ and a radial slide $d_3$ behind a tool of
length $L_t$,
\[
 d_3=\sqrt{x^2+y^2}-L_t,\qquad d_2=z,\qquad \theta_1=\operatorname{atan2}(\cdot,\cdot),
\]
the answer is unfinished until you add: $d_3\ge 0$, i.e.\ $\sqrt{x^2+y^2}\ge L_t$ --- points closer
to the column than the tool is long are unreachable --- and $z$ must lie inside the column's stroke.

How many branches should you expect? Count the sub-chains, not the joints. Elbow-up/elbow-down pairs
come from a \emph{two-revolute (2R) sub-chain}, where two link lengths and one target distance form a
triangle that can close in two mirror-image ways. A cylindrical arm has one revolute joint and no 2R
sub-chain, so its solution is \textbf{unique}. Reporting a $\pm$ branch for $d_3$ imports a pattern
from a different geometry: $d_3+L_t=\sqrt{x^2+y^2}$ is a length, so only the positive root is
physical.

Analytic IK uses robot geometry to isolate variables and enumerate branches. Numerical IK instead iteratively reduces pose error, commonly through a Jacobian-based update. Numerical methods generalise better but need an initial guess, a convergence test, joint-limit handling, and a strategy near singularities.
\source{\texttt{resources/lecture-01-manipulator-kinematics.pdf}, pp.~87--117.}

\paragraph{The tool-configuration vector: six numbers instead of twelve.}
Writing a tool pose as $(p,R)$ carries twelve numbers to express six degrees of freedom, and the
rotation matrix is awkward to use as the \emph{input} to an inverse-kinematics problem. The handout
therefore compresses orientation into the \emph{approach vector} $r_3$ --- the third column of $R$,
the direction the tool advances --- which fixes the tool's direction but says nothing about how far
it has rolled about that direction. Since $r_3$ is a unit vector, its \emph{length} is free, so the
roll can be hidden in the magnitude:
\[
 w=\begin{bmatrix}p\\[2pt] e^{q_n/\pi}\,r_3\end{bmatrix}\in\mathbb{R}^6 ,
\]
where $q_n$ is the \textbf{tool roll angle} --- the rotation \emph{about} the approach vector. The
exponential is chosen because $e^{q_n/\pi}>0$ always and is invertible, so the roll is recoverable:
$q_n=\pi\ln\lVert w_{4:6}\rVert$.

\paragraph{Common pitfall.}
$q_n$ is the roll \emph{about} $r_3$, not the first joint angle. Rotating the whole arm about its
base column swings the approach vector to a new \emph{direction}; it does not roll the tool about
it, so it does not enter the exponent. An arm with no roll degree of freedom --- a cylindrical
RPP arm, for instance --- has $q_n=0$, hence $e^{q_n/\pi}=1$ and $\lVert w_{4:6}\rVert=1$. If your
$w$ comes out with $\lVert w_{4:6}\rVert\ne 1$ on an arm that cannot roll, the decode
$q_n=\pi\ln\lVert w_{4:6}\rVert$ is claiming a roll the mechanism cannot perform.
\source{\texttt{resources/lecture-01-manipulator-kinematics.pdf}, pp.~91--95.}

\section{\enspace From a geometric path to a timed trajectory}
A path specifies where the tool travels; a trajectory also specifies when it occupies each point. Writing $w(t)=\Gamma(s(t))$ separates the geometric curve $\Gamma$ from the timing law $s(t)$. This distinction lets the same collision-free path be reused with slower or faster motion.

For one segment, a cubic polynomial
\[
 w(t)=at^3+bt^2+ct+d,qquad 0\leq t\leq T,
\]
can satisfy two endpoint positions and two endpoint velocities. Componentwise, for $w(0)=w_0$, $\dot w(0)=v_0$, $w(T)=w_1$, and $\dot w(T)=v_1$,
\begin{align*}
 d&=w_0,&c&=v_0,\\
 a&=\frac{2(w_0-w_1)+T(v_0+v_1)}{T^3},&
 b&=\frac{3(w_1-w_0)-T(2v_0+v_1)}{T^2}.
\end{align*}
The powers of $T$ are a dimensional check: $a$ has units of position/time$^3$ and $b$ position/time$^2$. Multi-segment motion should join with suitable velocity and acceleration continuity; matching positions alone can still create abrupt motion.

\paragraph{Building a grasp pose: which way do the jaws close?}
To write ${}^0T_{\text{pick}}$ for a parallel-jaw gripper you need two physical directions and then
one cross product; guessing all three columns is how orthonormality gets lost.
\begin{enumerate}
\item \textbf{Approach} $r_3$ --- the direction the tool \emph{advances into} the part. Picking
      from above means descending, so $r_3=-z_0=(0,0,-1)$, \emph{not} $+z_0$. The approach vector
      points from the tool toward the workpiece.
\item \textbf{Sliding} $r_2$ --- the direction the jaws \emph{close}. This is perpendicular to the
      faces they touch. ``Grasp the long sides'' names the two faces being contacted; those faces
      are separated \emph{across} the part, so the jaws travel along the \emph{short} axis. If the
      long axis lies along $x_0$, then $r_2=(0,1,0)$. Choosing $r_2$ along the long axis instead
      would drive the jaws at the part's short end faces.
\item \textbf{Normal} $r_1=r_2\times r_3$ --- computed, never guessed. Then verify
      $r_1\times r_2=r_3$ and $\det R=+1$. A right-handed orthonormal set can still be the
      \emph{wrong} frame for the task, so this check confirms validity, not correctness.
\item \textbf{Translation} --- the centroid of the part \emph{in its target location}. For a place
      pose that stacks part $A$ on part $B$, the height is
      $z = z_B + \tfrac12 h_B + \tfrac12 h_A$: rise to $B$'s top surface, then lift by $A$'s own
      half-height. Both half-heights, or the part ends up buried in the one below it. Three terms,
      every time --- and compute them from the dimensions in front of you, never from a number you
      remember from a similar problem.
\end{enumerate}

\paragraph{The pick and the place are two different orientations.}
Run the recipe above \emph{twice}, because $r_2$ depends on how the part is lying \emph{at that
moment}, and a place operation may well re-orient it:
\begin{itemize}
\item at the \textbf{pick}, use the part's \emph{current} long-axis direction;
\item at the \textbf{place}, use its \emph{target} long-axis direction --- which is whatever the task
      demands, for example ``aligned with the part underneath''.
\end{itemize}
If the two directions differ by $90^\circ$, so do the two gripper frames, and
$R_{\text{place}}=R_z(\pm90^\circ)R_{\text{pick}}$. If they are the same, $R$ is unchanged and only
the translation moves. Deciding which case you are in is the whole question --- and both frames are
perfectly valid rotations, so orthonormality and $\det R=+1$ will \emph{not} catch the mistake of
attaching the right matrix to the wrong pose.

The cheap habit: before writing either matrix, write the long-axis direction for that pose on its own
line (``long $\parallel y_0$'' then ``long $\parallel x_0$''). The jaw direction is then forced --- it
is the perpendicular one --- and the two poses cannot be interchanged by accident.
\source{\texttt{resources/lecture-01-manipulator-kinematics.pdf}, pp.~118--132; tool-frame column
naming (normal/sliding/approach) pp.~57--58.}

\section{\enspace Jacobians, velocity control, and singularities}
If the task coordinate is $x=w(q)$, differentiating with respect to time gives
\[
 \dot x=J(q)\dot q,
 \qquad J(q)=\frac{\partial w}{\partial q}.
\]
Column $j$ of $J$ is the instantaneous task-space motion produced by unit speed at joint $j$ while all other joints are fixed. This is a local linear relationship at the current configuration, not an exact formula for a large finite displacement.

When an exact inverse does not exist, the Moore--Penrose pseudoinverse gives a least-squares or minimum-norm rate command:
\[
 \dot q=J^\dagger\dot x.
\]
For a redundant robot, null-space motion $(I-J^\dagger J)z$ can change joint posture without changing the commanded first-order tool motion. This can help avoid limits or improve dexterity, but only if the secondary objective is chosen deliberately.

A singular configuration occurs when $J$ loses rank. Then at least one task-space direction becomes unattainable instantaneously, and a pseudoinverse command may demand very large joint rates near the singularity. The smallest singular value of $J$ is a useful warning: as it approaches zero, the mapping becomes ill-conditioned. In practice, slow down, re-plan, or use a damped least-squares inverse rather than trusting an unstable ordinary inverse.
\source{\texttt{resources/lecture-01-manipulator-kinematics.pdf}, pp.~133--156.}

\section{\enspace Full-course learning checklist}
\begin{itemize}
\item Can I distinguish a link, joint, end effector, pose, and workspace?
\item Can I explain why a reprogrammable robot differs from hard automation?
\item Can I state the source and target frame for every rotation or transformation?
\item Can I explain why the columns of a rotation matrix are axes expressed in another frame?
\item Can I invert a homogeneous transform without forgetting to rotate the translation?
\item Can I assign D--H frames, identify the variable parameter, and check the transform chain?
\item Before solving IK, can I check reachability, orientation freedom, limits, and solution branches?
\item Can I distinguish path geometry from trajectory timing and check cubic coefficient units?
\item Can I explain $\dot x=J\dot q$ and what rank loss means physically?
\end{itemize}
\end{document}
